% =====================================================================
%  04-implementacion.tex — Capitulo 4: Implementacion del prototipo
% =====================================================================

\chapter{Implementación del prototipo}\label{cap:implementacion}

Este capítulo documenta la implementación del prototipo cuyo análisis y diseño se presentó en el Capítulo~\ref{cap:analisis}. La implementación conserva la dirección de dependencias de ese diseño: el sistema se construyó como un workspace de Cargo de cinco crates con la regla de dependencias de la arquitectura hexagonal impuesta por el grafo de compilación (\cref{sec:workspace-cargo}), y cada uno de los servicios criptográficos especificados en el diseño criptográfico (\cref{sec:diseno-cripto}) se materializó mediante puertos de dominio, casos de uso y adaptadores concretos de infraestructura, tal como anticipó la definición de la arquitectura hexagonal (\cref{sec:hexagonal-definicion}). El capítulo se organiza en seis secciones. La primera describe la estructura real del workspace y el andamiaje de calidad que lo rodea: perfil de compilación, configuración, telemetría, integración continua y política de dependencias. La segunda recorre los diez servicios criptográficos implementados, identificando para cada uno su puerto, su adaptador, los crates empleados y las decisiones técnicas registradas. La tercera presenta la interfaz de línea de comandos que expone los casos de uso. La cuarta documenta la aplicación HTTP autenticada, PostgreSQL, Redis y la persistencia documental cifrada. La quinta reporta los resultados de la exploración del sellado de tiempo, el componente que el análisis de riesgos señaló como el de mayor exposición (\cref{sec:riesgos-tecnicos}). La sexta registra las divergencias respecto al diseño del Capítulo~\ref{cap:analisis} y las decisiones de arquitectura documentadas durante la construcción.

\section{Estructura del workspace y andamiaje}\label{sec:impl-workspace}

\textbf{Estado implementado al 12 de septiembre de 2026.} El prototipo expone una API Axum multiusuario con alta inicial, autenticación Argon2id, TOTP o recuperación, sesiones revocables, expedientes y asignaciones. Las operaciones documentales exigen un expediente y comprueban el acceso tanto antes de preparar el resultado como al confirmarlo. PostgreSQL persiste usuarios, expedientes, pertenencias, documentos cifrados y una cadena de auditoría compartida; Redis conserva el estado efímero de autenticación. Están implementados carga, sellado local, verificación y exportación de evidencia. Permanecen pendientes la interfaz de producto, las consultas e historial de versiones documentales y la actividad procesal.

{\emergencystretch=2em El prototipo es un workspace de Cargo con cinco crates miembros bajo \rutarepo{crates/}: \texttt{domain}, \texttt{application}, \texttt{infrastructure}, \texttt{web} y \texttt{bin}, conforme a la estructura de \cref{sec:workspace-cargo}. La dirección de dependencias apunta al dominio: \texttt{domain} no depende de otro crate del workspace; \texttt{application} depende de \texttt{domain}; \texttt{infrastructure} y \texttt{web} dependen de las dos capas interiores; y \texttt{bin}, raíz de composición que produce \texttt{despacho-cli}, depende de todas. Este último es el único crate que utiliza \texttt{anyhow} en la frontera del proceso. HTTP consume los puertos de aplicación \texttt{Identity\allowbreak Workflow}, \texttt{Case\allowbreak Workflow} y \texttt{Case\allowbreak Document\allowbreak Workflow}, sin conocer los adaptadores concretos. La \cref{tab:impl-crates} sintetiza las responsabilidades; las pruebas unitarias viven en \texttt{src} y las de integración en \texttt{tests/}. La convención de mantenimiento limita a menos de 400 líneas cada archivo de código, según \texttt{CONTRIBUTING.md}.\par}

\begin{table}[H]
  \centering
  \caption{Responsabilidades de los crates del workspace.}
  \label{tab:impl-crates}
  \begin{tabularx}{\linewidth}{@{}l X X@{}}
    \toprule
    \textbf{Crate} & \textbf{Responsabilidad} & \textbf{Elementos del flujo integrado} \\
    \midrule
    \texttt{domain} & Valores, reglas, puertos y errores; sin E/S & Identidad, expedientes, permisos, criptografía y cadena de auditoría. \\
    \texttt{application} & Casos de uso y fronteras de persistencia & Autenticación, asignaciones y operaciones documentales por expediente. \\
    \texttt{infrastructure} & Adaptadores de servicios y almacenamiento & PostgreSQL, Redis, cifrado, firma, TSA, ZIP e importación local. \\
    \texttt{web} & Adaptador HTTP y DTO & Extracción de peticiones, rutas, errores y límites de concurrencia. \\
    \texttt{bin} & Composición y comandos & Servidor, preparación de esquema, importación y herramientas offline. \\
    \bottomrule
  \end{tabularx}
\end{table}

La regla de dependencias no descansa en la disciplina del equipo sino en dos mecanismos verificables. El primero es el propio grafo de Cargo: como se argumentó en el diseño (\cref{sec:workspace-cargo}), una importación prohibida produce un error de compilación, no una advertencia de revisión. El segundo es una prueba de arquitectura, \rutarepo{crates/bin/tests/architecture.rs}, que lee el manifiesto del crate \texttt{domain} y falla si este declara dependencia alguna de \texttt{application}, \texttt{infrastructure} o \texttt{web}; esta prueba protege contra la degradación gradual del manifiesto, que el compilador no reportaría mientras los símbolos prohibidos no se usen. Una segunda prueba estructural, \rutarepo{crates/bin/tests/port_coverage.rs}, recorre el workspace y exige que cada \texttt{pub trait} del dominio ---cada puerto--- tenga al menos un adaptador (\texttt{impl <Trait> for}) en \texttt{infrastructure} y al menos un mock generado con \texttt{mockall} en los tests; para detectar un escáner degradado, la prueba usa puertos centinela (\texttt{AuditLog}, \texttt{Document\allowbreak Hasher}, \texttt{KeyManager}) cuya ausencia delataría un análisis roto.

La versión mínima soportada del compilador es Rust 1.88, declarada mediante \texttt{rust-version} en el \texttt{Cargo.toml} raíz y comprobada por un job obligatorio de integración continua. La elevación desde 1.78 permite adoptar correcciones de seguridad de \texttt{time} y de la familia PostgreSQL, conforme a \rutarepo{docs/adr/0013-raise-msrv-for-security-fixes.md}. Se conserva la familia de APIs construida alrededor de \texttt{der}~0.7 ---\texttt{x509-cert}~0.2, \texttt{cms}~0.2, \texttt{x509-tsp}~0.1 y \texttt{rsa}~0.9--- porque su sustitución exige una migración coordinada, no porque 1.78 siga siendo el requisito vigente. Las dependencias compartidas se declaran en \texttt{[workspace.dependencies]}. El perfil de release aplica \texttt{strip = true} y \texttt{lto = true}; en desarrollo se optimizan los núcleos \texttt{argon2}, \texttt{blake2}, \texttt{num-bigint-dig} y \texttt{rsa}, manteniendo información de depuración en el resto del workspace.

La configuración se carga desde variables de entorno con \texttt{dotenvy} (archivo \texttt{.env} local, nunca versionado; la plantilla comentada \texttt{.env.example} sí lo está). La \cref{tab:impl-env} lista las variables. Dos decisiones de manejo de secretos son transversales: la carga de configuración registra en la bitácora únicamente la \emph{presencia} de cada variable, jamás su valor; y todo material secreto leído del entorno o de la entrada estándar (la KEK, la credencial del proveedor de sellado, las contraseñas) aterriza en buffers \texttt{Zeroizing} que sobrescriben su contenido al liberarse, conforme al ADR \rutarepo{docs/adr/0003-zeroize-secret-material-in-domain.md}.

\begin{table}[H]
  \centering
  \caption{Variables de entorno del prototipo (plantilla \texttt{.env.example}).}
  \label{tab:impl-env}
  \begin{tabularx}{\linewidth}{@{}l X@{}}
    \toprule
    \textbf{Variable} & \textbf{Propósito} \\
    \midrule
    \texttt{RUST\_LOG} & Filtro de severidad de la telemetría (por defecto \texttt{info}). \\
    \texttt{CINCEL\_BASE\_URL} & URL base del proveedor remoto de sellado de tiempo. \\
    \texttt{CINCEL\_API\_KEY} & Credencial del proveedor de sellado (secreto; se llena localmente). \\
    \texttt{DATABASE\_URL} & PostgreSQL para usuarios, expedientes, documentos y auditoría; conexión administrativa para migrar y rol restringido para servir. \\
    \texttt{REDIS\_URL} & Cadena de conexión Redis para sesiones, desafíos y límites. \\
    \texttt{KEK\_BASE64} & Clave de cifrado de claves: 32 bytes aleatorios en base64 (secreto). \\
    \texttt{NEW\_KEK\_BASE64} & KEK de reemplazo; la lee únicamente \texttt{vault rotate-kek}. \\
    \texttt{AUDIT\_LOG\_PATH} & Ruta de la bitácora offline; HTTP usa PostgreSQL (por defecto \texttt{audit-log.jsonl}). \\
    \texttt{PKI\_CA\_DIR} & Directorio de trabajo de la CA interna (por defecto \texttt{pki-ca}). \\
    \texttt{TSA\_DIR} & Directorio de trabajo de la TSA local (por defecto \texttt{pki-tsa}). \\
    \bottomrule
  \end{tabularx}
\end{table}

La telemetría usa \texttt{tracing} y \texttt{tracing-subscriber} con filtro por entorno, materializando la observabilidad estructurada prevista en el diseño (\cref{sec:observabilidad}). Todos los diagnósticos se emiten por la salida de error estándar, de modo que la salida estándar queda limpia para resultados y tuberías de shell. La promesa de que ningún secreto aparece en los diagnósticos no es una convención: la prueba \rutarepo{crates/bin/tests/log_redaction.rs} ejecuta los flujos con el nivel de log más verboso e inspecciona lo emitido, fallando si algún material secreto conocido aparece en la salida de error.

El workflow de integración continua ejecuta siete jobs en cada push y pull request (\cref{tab:impl-ci}). La comprobación de MSRV es obligatoria y usa Rust 1.88. Los jobs de pruebas y cobertura preparan PostgreSQL y Redis y separan las bases de identidad, expedientes y documentos, de modo que las pruebas de persistencia no se omiten por ausencia de configuración. El umbral de \rutarepo{scripts/coverage-gate.sh} exige 90\,\% de líneas por crate en \texttt{domain}, \texttt{application} e \texttt{infrastructure}; \texttt{bin} se reporta sin bloquear. Estos valores son criterios de aceptación configurados, no mediciones de una ejecución concreta.

\begin{table}[H]
  \centering
  \caption{Jobs del workflow de integración continua.}
  \label{tab:impl-ci}
  \begin{tabularx}{\linewidth}{@{}l l X@{}}
    \toprule
    \textbf{Job} & \textbf{Herramienta} & \textbf{Criterio de aprobación} \\
    \midrule
    format & \texttt{cargo fmt} & Formato canónico en todo el workspace. \\
    lint & \texttt{cargo clippy} & Cero advertencias (\texttt{-D warnings}). \\
    test & \texttt{cargo test} & Suite completa del workspace en verde. \\
    msrv & \texttt{cargo check} & Compila con Rust 1.88; comprobación obligatoria. \\
    coverage & \texttt{cargo-llvm-cov} & $\geq$ 90\,\% de líneas por crate criptográfico. \\
    deny & \texttt{cargo-deny} & Avisos, licencias, duplicados y fuentes en regla. \\
    binary-size & \texttt{cargo build -{}-release} & \texttt{despacho-cli} $\leq$ 25~MB. \\
    \bottomrule
  \end{tabularx}
\end{table}

El diseño comprometió cinco fitness functions como verificaciones automatizadas de los atributos de calidad (\cref{sec:fitness-functions}) \cite{ford2023}. La \cref{tab:impl-fitness} muestra dónde vive cada una en el repositorio. Cuatro de las cinco diseñadas se implementaron directamente; la función que exige un span de \texttt{tracing} por handler HTTP sigue pendiente aun cuando los handlers ya existen, y en su lugar se incorporó el umbral de cobertura por crate como función de aptitud adicional.

\begin{table}[H]
  \centering
  \caption{Fitness functions implementadas y su ubicación en el repositorio.}
  \label{tab:impl-fitness}
  \begin{tabularx}{\linewidth}{@{}l L{6.4cm} X@{}}
    \toprule
    \textbf{Id} & \textbf{Propiedad verificada} & \textbf{Verificación implementada} \\
    \midrule
    FF-1 & El dominio no depende de capas externas & Grafo de Cargo en cada build y \rutarepo{crates/bin/tests/architecture.rs}. \\
    FF-2 & Ningún log contiene material secreto & \rutarepo{crates/bin/tests/log_redaction.rs} al nivel de log más verboso. \\
    FF-3 & Todo puerto tiene adaptador y mock & \rutarepo{crates/bin/tests/port_coverage.rs} con puertos centinela. \\
    FF-4 & Cobertura $\geq$ 90\,\% en crates criptográficos & \rutarepo{scripts/coverage-gate.sh} sobre el reporte de \texttt{cargo-llvm-cov}. \\
    FF-5 & Binario release $\leq$ 25~MB & Job \texttt{binary-size} del workflow de integración continua. \\
    \bottomrule
  \end{tabularx}
\end{table}

La política de dependencias revisa avisos, licencias, duplicados y fuentes del grafo. Se configura en \texttt{deny.toml} para \texttt{cargo-deny} y en \texttt{audit.toml} para \texttt{cargo-audit}. Permanece declarada la excepción RUSTSEC-2023-0071 \cite{rustsec2023_0071}, relativa a temporización en operaciones de clave privada del crate \texttt{rsa}. El razonamiento original de \rutarepo{docs/adr/0002-rsa-signing-crate-and-advisory.md} suponía firma por CLI; esa premisa no describe la exposición HTTP actual. La autenticación y los límites de concurrencia acotan quién puede solicitar firmas y cuántas se ejecutan, pero no corrigen el canal lateral. Antes de un despliegue público se requiere revisar el modelo de amenaza y resolver la elección del firmante. La API local no se presenta como una configuración de producción validada.

\section{Servicios criptográficos}\label{sec:impl-servicios}

El diseño criptográfico del Capítulo~\ref{cap:analisis} (\cref{sec:diseno-cripto}) especificó los servicios de integridad, confidencialidad, no repudio y autenticación del prototipo. Esta sección los recorre como diez servicios implementados, S1 a S10, cada uno con la misma anatomía: un puerto (trait) en \texttt{domain}, un adaptador en \texttt{infrastructure} y un caso de uso en \texttt{application} que los orquesta. La \cref{tab:impl-servicios} presenta el mapa completo. Dos convenciones atraviesan todos los servicios. Primera: los rechazos criptográficos son \emph{resultados} tipados (enumeraciones como \texttt{Signature\allowbreak Verification} o \texttt{Certificate\allowbreak Validation}), nunca errores; el camino de error queda reservado a fallos del backend o a evidencia inevaluable, lo que permite construir políticas futuras ---por ejemplo, bloqueo de cuenta--- sin confundir un intento fallido con una falla del sistema. Segunda: la delegación a la herramienta \texttt{openssl} como subproceso ocurre en exactamente tres puntos ---la CA interna, la TSA local y la etapa criptográfica del verificador RFC 3161---; todo lo demás es Rust nativo sobre \texttt{ring}, \texttt{rsa}, \texttt{sha2}, \texttt{x509-cert}, \texttt{cms}, \texttt{x509-tsp}, \texttt{argon2} y \texttt{totp-rs}.

El manejo de errores materializa el patrón de errores tipados por capa que el diseño estableció (\cref{sec:manejo-errores}). El crate de dominio define \texttt{DomainError} (\rutarepo{crates/domain/src/error.rs}) con 26 variantes tipadas mediante \texttt{thiserror}, una por condición de fallo con significado propio ---longitud de digest inválida, payload sellado malformado, CRL vencida, nombre de entrada de archivo inutilizable---. Una variante merece mención explícita: \texttt{Authentication\allowbreak Failed}, el rechazo de un descifrado autenticado, es deliberadamente opaca y no revela si falló la clave, el tag o los datos, negando a un adversario la retroalimentación que distingue esos casos. Los errores de los backends de infraestructura (\texttt{CryptoError}, \texttt{TsaError}) se convierten hacia el dominio mediante conversiones \texttt{From}, y la capa de aplicación añade en \texttt{Application\allowbreak Error} solo lo que le es propio: el fallo genérico de un puerto, un archivo vault malformado y un certificado recién emitido que no supera su validación posterior a la emisión.

\begin{table}[H]
  \centering
  \caption{Servicios criptográficos implementados: puertos del dominio y adaptadores de infraestructura.}
  \label{tab:impl-servicios}
  \begin{tabularx}{\linewidth}{@{}l L{3.4cm} L{4.6cm} X@{}}
    \toprule
    \textbf{Id} & \textbf{Servicio} & \textbf{Puerto(s) del dominio} & \textbf{Adaptador(es)} \\
    \midrule
    S1 & Hash de documentos & \texttt{Document\allowbreak Hasher} & \texttt{Ring\allowbreak Sha256\allowbreak Hasher} \\
    S2 & Cifrado autenticado y envelope encryption & \texttt{Authenticated\allowbreak Cipher}, \texttt{KeyManager} & \texttt{Ring\allowbreak Aes\allowbreak Gcm\allowbreak Cipher}, \texttt{Envelope\allowbreak Key\allowbreak Manager} \\
    S3 & PKI interna & \texttt{Certificate\allowbreak Authority}, \texttt{Certificate\allowbreak Validator} & \texttt{Openssl\allowbreak Ca\allowbreak Adapter}, \texttt{X509\allowbreak Chain\allowbreak Validator} \\
    S4 & Firma digital & \texttt{Document\allowbreak Signer}, \texttt{Signature\allowbreak Verifier} & \texttt{Rsa\allowbreak Pkcs1\allowbreak Signer}, \texttt{Rsa\allowbreak Pkcs1\allowbreak Verifier} \\
    S5 & Sellado de tiempo & \texttt{Timestamp\allowbreak Service}, \texttt{Timestamp\allowbreak Verifier} & \texttt{Local\allowbreak Openssl\allowbreak Tsa}, \texttt{Cincel\allowbreak Tsa\allowbreak Adapter}, \texttt{Rfc3161\allowbreak Verifier} \\
    S6 & Verificación integral & (compone los de S1, S3, S4 y S5) & caso de uso \texttt{Verify\allowbreak Document} \\
    S7 & Hash de contraseñas & \texttt{Password\allowbreak Hasher} & \texttt{Argon2id\allowbreak Hasher} \\
    S8 & Segundo factor TOTP & \texttt{Totp\allowbreak Provider}, \texttt{Recovery\allowbreak Code\allowbreak Generator} & \texttt{Totp\allowbreak Rs\allowbreak Provider}, \texttt{Random\allowbreak Recovery\allowbreak Code\allowbreak Generator} \\
    S9 & Cadena de auditoría & \texttt{AuditLog}, \texttt{Clock} & \texttt{Postgres\allowbreak Audit\allowbreak Log}, \texttt{File\allowbreak Audit\allowbreak Log}, \texttt{SystemClock} \\
    S10 & Paquete de evidencia & \texttt{Archive\allowbreak Writer} & \texttt{Stored\allowbreak Zip\allowbreak Writer} \\
    \bottomrule
  \end{tabularx}
\end{table}

\subsection{S1. Hash de documentos con SHA-256}\label{sec:impl-s1-hash}

El puerto \texttt{Document\allowbreak Hasher} (\rutarepo{crates/domain/src/crypto/hasher.rs}) es el más simple del sistema y sirve para ilustrar el patrón que repiten todos los demás: un trait del dominio que declara el contrato criptográfico sin nombrar backend alguno. El Listado~\ref{lst:impl-puerto-hasher} lo reproduce íntegro.

\begin{lstlisting}[language=Rust,caption={Puerto de hashing del dominio (\texttt{crates/domain/src/crypto/hasher.rs}).},label={lst:impl-puerto-hasher}]
/// Computes SHA-256 digests of document content.
///
/// Adapters implement this port with a concrete cryptographic backend. Both
/// operations must produce the same digest for the same input bytes.
pub trait DocumentHasher {
    /// Hashes a byte slice already held in memory.
    fn hash_bytes(&self, data: &[u8]) -> Sha256Digest;

    /// Hashes a stream by consuming it in fixed-size blocks, so inputs of
    /// any length are digested without loading them fully into memory.
    ///
    /// # Errors
    ///
    /// Returns [`DomainError::StreamRead`] when reading from `reader` fails.
    fn hash_stream(&self, reader: &mut dyn Read) -> Result<Sha256Digest, DomainError>;
}
\end{lstlisting}

El adaptador \texttt{Ring\allowbreak Sha256\allowbreak Hasher} (\rutarepo{crates/infrastructure/src/hashing.rs}) lo implementa con el crate \texttt{ring}~0.17. La variante de flujo consume la entrada en bloques de 64~KiB, de modo que el consumo de memoria es plano para entradas de cualquier tamaño, y reintenta las lecturas interrumpidas por señales del sistema operativo. El resultado es el value object \texttt{Sha256\allowbreak Digest}, un arreglo fijo de 32 bytes con constructores que rechazan longitudes o codificaciones hexadecimales inválidas. Siguiendo la práctica de desarrollo guiado por pruebas del proyecto, los primeros tests del adaptador son los vectores oficiales de FIPS 180-4 \cite{nist2015fips1804}, complementados con una verificación de interoperabilidad contra \texttt{sha256sum} y con propiedades verificadas por generación aleatoria de casos. El caso de uso \texttt{Hash\allowbreak Document} expone el servicio a la interfaz de usuario.

\subsection{S2. Cifrado autenticado y envelope encryption}\label{sec:impl-s2-cifrado}

{\emergencystretch=2em La confidencialidad en reposo diseñada en \cref{sec:cifrado-reposo} se implementa con dos puertos. El primero, \texttt{Authenticated\allowbreak Cipher} (\rutarepo{crates/domain/src/crypto/cipher.rs}), declara las operaciones de sellado y apertura autenticados; el material de clave cruza la interfaz como slices de bytes para que el crate de dominio permanezca libre de dependencias criptográficas, mientras los buffers propietarios viven en contenedores \texttt{Zeroizing} del lado del llamador. Su adaptador, \texttt{Ring\allowbreak Aes\allowbreak Gcm\allowbreak Cipher} (\rutarepo{crates/infrastructure/src/encryption.rs}), implementa AES-256-GCM con \texttt{ring} \cite{dworkin2007,mcgrew2008}: en cada operación de sellado se extrae un nonce fresco de 96 bits del CSPRNG del sistema, y la salida adopta el formato fijo nonce~(12~bytes)~$\Vert$~criptograma~$\Vert$~tag~(16~bytes), representado por el value object \texttt{Sealed\allowbreak Payload}, que rechaza cualquier secuencia menor a 28 bytes. Todo rechazo de descifrado se colapsa deliberadamente en el error opaco \texttt{Authentication\allowbreak Failed}, que no revela si falló la clave, el tag o los datos.\par}

Los datos adicionales autenticados ligan cada criptograma a la identidad y versión del documento, tal como exigió el diseño para prevenir ataques de intercambio y de reversión. La función de dominio \texttt{document\_aad} (Listado~\ref{lst:impl-document-aad}) construye ese vínculo: 16 bytes crudos del UUID del documento seguidos de la versión como 4 bytes big-endian.

\begin{lstlisting}[language=Rust,caption={Construcción del AAD que liga el criptograma a su documento (\texttt{crates/domain/src/crypto/cipher.rs}).},label={lst:impl-document-aad}]
/// Builds the additional authenticated data that binds a sealed document to
/// its identity: the 16 raw bytes of the document id UUID followed by the
/// document version as 4 big-endian bytes.
///
/// Sealing a document with this AAD makes decryption fail unless the caller
/// asks for exactly the same document and version. That prevents
/// substituting one document's ciphertext for another's (the id would
/// differ) and replaying an older version of the same document (the version
/// would differ).
pub fn document_aad(id: DocumentId, version: DocumentVersion) -> [u8; DOCUMENT_AAD_LEN] {
    let mut aad = [0u8; DOCUMENT_AAD_LEN];
    aad[..16].copy_from_slice(id.as_uuid().as_bytes());
    aad[16..].copy_from_slice(&version.get().to_be_bytes());
    aad
}
\end{lstlisting}

El segundo puerto, \texttt{KeyManager} (\rutarepo{crates/domain/src/crypto/keys.rs}), gobierna el esquema de envelope encryption \cite{dworkin2012}: generar una DEK, envolverla bajo la KEK, desenvolverla y re-envolverla bajo una KEK nueva. Su adaptador \texttt{Envelope\allowbreak Key\allowbreak Manager} genera una DEK aleatoria de 32 bytes por documento y la sella con el mismo cifrador AES-256-GCM, pero con AAD \emph{vacío}: la decisión, documentada en el propio código, es que la DEK envuelta ya queda ligada a su documento por el archivo que la transporta, y mantener la envoltura independiente de la identidad permite rotar la KEK re-envolviendo todas las claves sin necesidad de saber a qué documento pertenece cada una. El caso de uso \texttt{Encrypt\allowbreak Document} serializa el resultado en el formato de archivo \texttt{DVLT1} (\cref{tab:impl-vault}), definido normativamente en \rutarepo{crates/application/src/vault.rs}; \texttt{Decrypt\allowbreak Document} reconstruye el AAD con la identidad y versión que el llamador afirma, de modo que un identificador o versión incorrectos fallan exactamente igual que un criptograma alterado; y \texttt{RotateKek} re-envuelve únicamente la DEK, copiando intactos los bytes del documento sellado, con un costo independiente del tamaño del documento ---la propiedad de rotación que el diseño exigió en \cref{sec:cifrado-reposo}.

\begin{table}[H]
  \centering
  \caption[Formato del archivo vault \texttt{DVLT1}]{Formato del archivo vault \texttt{DVLT1} (definición normativa en \texttt{crates/application/src/vault.rs}).}
  \label{tab:impl-vault}
  \begin{tabularx}{\linewidth}{@{}l l X@{}}
    \toprule
    \textbf{Desplazamiento} & \textbf{Tamaño} & \textbf{Campo} \\
    \midrule
    0 & 5 & Bytes mágicos ASCII \texttt{DVLT1}. \\
    5 & 4 & Longitud $N$ (big-endian) de la DEK envuelta. \\
    9 & $N$ & DEK envuelta: nonce~$\Vert$~criptograma~$\Vert$~tag, sellada bajo la KEK con AAD vacío. \\
    $9+N$ & resto & Documento sellado: nonce~$\Vert$~criptograma~$\Vert$~tag, sellado bajo la DEK con AAD = identificador del documento (16 bytes UUID)~$\Vert$~versión (4 bytes big-endian). \\
    \bottomrule
  \end{tabularx}
\end{table}

\subsection{S3. Infraestructura de clave pública interna}\label{sec:impl-s3-pki}

La CA interna de una sola capa especificada en \cref{sec:firma-digital} se implementó como un conjunto de scripts versionados bajo \rutarepo{pki/} ---\texttt{init-ca.sh}, \texttt{issue-cert.sh}, \texttt{revoke.sh}, \texttt{gen-crl.sh}--- que operan la herramienta OpenSSL con configuración declarativa (\texttt{openssl.cnf}). Los parámetros coinciden con el diseño: claves RSA de 3\,072 bits, firma sha256WithRSAEncryption en todos los casos, certificado raíz autofirmado con vigencia de 1\,825 días (cinco años), certificados de entidad final con vigencia de 365 días y CRL con vigencia de siete días conforme a RFC 5280 \cite{cooper2008}. La política de emisión obliga a que país y organización del certificado emitido coincidan con los de la CA, ignora las extensiones solicitadas en el CSR (\texttt{copy\_extensions = none}) y aplica un retroceso deliberado de cinco minutos al inicio de vigencia ---práctica común de las autoridades públicas--- para que un certificado sea utilizable de inmediato aun con desfase de reloj entre máquinas. El puerto \texttt{Certificate\allowbreak Authority} se implementa en \texttt{Openssl\allowbreak Ca\allowbreak Adapter} (\rutarepo{crates/infrastructure/src/certificates/authority.rs}), que ejecuta esos scripts como subprocesos con el directorio de trabajo inyectado vía \texttt{PKI\_CA\_DIR} y normaliza los números de serie a hexadecimal en mayúsculas antes de invocar subproceso alguno. El caso de uso \texttt{Issue\allowbreak Certificate} rehúsa reportar éxito hasta validar el certificado recién emitido contra la raíz, de modo que una CA mal inicializada falla en la emisión y no en el primer uso. Los scripts operan con disciplina defensiva ---\texttt{set -euo pipefail}, verificación previa de que \texttt{openssl} está disponible, inicialización idempotente que rehúsa sobrescribir una CA existente--- y el material de claves privadas jamás se versiona: vive únicamente en el directorio de trabajo en tiempo de ejecución, con el \texttt{.gitignore} del repositorio como red de seguridad adicional; el \texttt{README.md} del directorio documenta el orden de uso, los comandos de verificación independiente y la versión de OpenSSL (3.5.7) con la que fueron probados.

La validación es el puerto \texttt{Certificate\allowbreak Validator}, implementado por \texttt{X509\allowbreak Chain\allowbreak Validator} (\rutarepo{crates/infrastructure/src/certificates/validator.rs}). Las crates \texttt{x509-cert}~0.2 y \texttt{der}~0.7 solo \emph{parsean} X.509 y deliberadamente no verifican firmas, así que la verificación criptográfica se implementó re-codificando a DER la estructura to-be-signed parseada y comprobando la firma PKCS\#1 v1.5 con \texttt{rsa} y \texttt{sha2} contra la clave pública del emisor ---la misma técnica para la firma de la CRL sobre su \texttt{tbsCertList}---, decisión registrada en \rutarepo{docs/adr/0004-certificate-validation-implementation.md} y correcta porque DER es una codificación canónica: decodificar y re-codificar entrada DER válida reproduce los bytes firmados. El validador ejecuta sus comprobaciones en orden fijo y reporta exactamente el primer fallo: (1) confianza del emisor ---nombre, algoritmo y firma---; (2) ventana de vigencia del certificado; (3) ventana de vigencia \emph{del emisor}, porque un emisor fuera de su propia vigencia no ancla confianza en ese instante; y (4) revocación contra la CRL opcional. El contrato de revocación es estricto: si se presenta una CRL, cualquier defecto de la lista ---ilegible, no firmada por el emisor, o vencida en su \texttt{nextUpdate}--- es un error duro y nunca un pase permisivo; omitir la CRL es una decisión explícita del llamador que el reporte de verificación hace visible. El instante de evaluación siempre es un timestamp explícito, lo que permite a las pruebas controlar el reloj.

\subsection{S4. Firma digital RSA-3072}\label{sec:impl-s4-firma}

El esquema de firma es el diseñado en \cref{sec:firma-digital}: RSA-3072 con SHA-256 y padding PKCS\#1 v1.5. Los puertos son \texttt{Document\allowbreak Signer} y \texttt{Signature\allowbreak Verifier} (\rutarepo{crates/domain/src/crypto/signature.rs}); nunca se firma el documento sino su digest SHA-256, calculado por el puerto de hashing, y el rechazo de una verificación es un resultado tipado, no un error. Los adaptadores \texttt{Rsa\allowbreak Pkcs1\allowbreak Signer} y \texttt{Rsa\allowbreak Pkcs1\allowbreak Verifier} (\rutarepo{crates/infrastructure/src/signing.rs}) usan los crates \texttt{rsa}~0.9 y \texttt{sha2}~0.10 con la feature \texttt{oid}, que liga SHA-256 a su identificador en la estructura DigestInfo del padding, de modo que los 32 bytes ya calculados se firman sin hashear dos veces. El firmador rechaza en construcción cualquier módulo menor a 3\,072 bits; acepta la clave privada en PEM PKCS\#8 o PKCS\#1, la retiene en un buffer \texttt{Zeroizing}, la parsea únicamente durante la operación de firma y la descarta al salir del ámbito, y su representación de depuración omite el material (\texttt{private key pem withheld}). El verificador acepta como material de clave tanto un certificado X.509 ---del que extrae el SubjectPublicKeyInfo--- como una clave pública SPKI en PEM; una firma cuya longitud no coincide con el tamaño de la clave se rechaza como malformada, y cualquier comprobación fallida se colapsa en \texttt{Mismatched\allowbreak Document\allowbreak Or\allowbreak Key}, porque un fallo de verificación no distingue un documento alterado de una clave ajena. La elección del crate \texttt{rsa} ---con el aviso RUSTSEC-2023-0071 aceptado y un plan de contingencia hacia \texttt{ring} si el riesgo se volviera inaceptable--- está registrada en \rutarepo{docs/adr/0002-rsa-signing-crate-and-advisory.md}.

\subsection{S5. Sellado de tiempo RFC 3161}\label{sec:impl-s5-sellado}

El sellado de tiempo conserva el contrato del diseño (\cref{sec:sellado-tiempo}): el puerto \texttt{Timestamp\allowbreak Service} recibe un digest SHA-256 y devuelve el token como bytes DER opacos, y el puerto \texttt{Timestamp\allowbreak Verifier} verifica un token contra el digest esperado y un ancla de confianza en PEM \cite{adams2001}. Existen dos adaptadores emisores, intercambiables sin tocar el dominio. \texttt{Local\allowbreak Openssl\allowbreak Tsa} (\rutarepo{crates/infrastructure/src/timestamp/local.rs}) es una autoridad de sellado local de desarrollo y demostración que ejecuta \texttt{openssl ts -query} y \texttt{openssl ts -reply} como subprocesos contra el directorio de trabajo preparado por \rutarepo{pki/issue-tsa-cert.sh}; el paso de respuesta se serializa entre hilos y procesos con un lock advisory exclusivo sobre el archivo \texttt{serial.lock} (crate \texttt{fs2}), porque OpenSSL avanza el archivo de serie con una secuencia leer-incrementar-escribir no atómica que, sin serialización, acuñaría números de serie duplicados. \texttt{Cincel\allowbreak Tsa\allowbreak Adapter} (\rutarepo{crates/infrastructure/src/timestamp/cincel.rs}) es el camino remoto hacia el sandbox del PSC \cite{cincel2024}, implementado con \texttt{reqwest}~0.12 en modo bloqueante sobre \texttt{rustls}; envía el digest por \texttt{POST} a \rutarepo{/api/v1/timestamps} con la credencial en el header \texttt{X-Api-Key}, maneja las respuestas de token inmediato, procesamiento diferido con sondeo (\texttt{RetryPolicy} configurable) y rechazo, y mantiene la credencial en un buffer \texttt{Zeroizing} excluido de la depuración y de todo mensaje de error. Los resultados de la exploración con ambos emisores se reportan en \cref{sec:impl-spike-tsa}.

{\emergencystretch=2em La verificación, \texttt{Rfc3161\allowbreak Verifier} (\rutarepo{crates/infrastructure/src/timestamp/verify.rs}), opera en dos etapas conforme a \rutarepo{docs/adr/0005-rfc3161-verification-strategy.md}. La etapa nativa decodifica el DER con \texttt{cms}~0.2 y \texttt{x509-tsp}~0.1 ---aceptando una respuesta \texttt{Time\allowbreak Stamp\allowbreak Resp} completa o un token CMS desnudo---, extrae el \texttt{TSTInfo}, compara el message imprint contra el digest esperado exigiendo tanto el OID id-sha256 (2.16.840.1.101.3.4.2.1) como la igualdad de bytes, y lee el tiempo de generación en formato RFC 3339; las entradas indecodificables y los desajustes de imprint se diagnostican aquí, con causas precisas y sin lanzar subproceso alguno. La etapa delegada encarga la comprobación de la firma CMS y de la cadena de certificados a \texttt{openssl ts -verify}, montando los archivos en un directorio temporal privado (modo 0700, nombre impredecible) que se elimina en todo camino de retorno; la aceptación produce \texttt{Valid} con el instante de generación, y un rechazo produce \texttt{Untrusted\allowbreak Token} con el diagnóstico de la herramienta. Implementar a mano la verificación CMS ---atributos firmados, comprobaciones de tipo de contenido, construcción de cadena--- habría añadido una superficie criptográfica sustancial que el prototipo no debe poseer, mientras que la dependencia de la herramienta \texttt{openssl} ya existía para la CA interna.\par}

\subsection{S6. Verificación integral}\label{sec:impl-s6-verificacion}

{\emergencystretch=2em El caso de uso \texttt{Verify\allowbreak Document} (\rutarepo{crates/application/src/verification.rs}) compone cuatro puertos ---\texttt{Document\allowbreak Hasher}, \texttt{Signature\allowbreak Verifier}, \texttt{Certificate\allowbreak Validator} y \texttt{Timestamp\allowbreak Verifier}--- en un único reporte sobre un flujo de entrada, materializando el flujo de verificación en cuatro pasos que el diseño describió en \cref{sec:sellado-tiempo}. El reporte contiene cuatro componentes (integridad, firma, certificado y sello de tiempo), cada uno con estado \texttt{Passed}, \texttt{Failed} o \texttt{Skipped} y una causa legible. La semántica es deliberadamente conservadora: los componentes cuya evidencia no se aportó se reportan omitidos \emph{con su razón}, nunca desaparecen en silencio, y no cuentan contra el veredicto; el veredicto es \texttt{NotValid} si cualquier componente ejercitado falló (Listado~\ref{lst:impl-verdict}). La integridad compara el digest recomputado contra el que atestigua el token de sellado ---sin token no existe referencia independiente y el componente se omite---, y el verificador revisa el imprint \emph{antes} que la cadena de confianza, de modo que un token no confiable aún puede confirmar que el digest coincide (integridad aprobada con sello fallido). El estado del certificado habla únicamente del instante de evaluación: que esté hoy expirado o revocado no invalida por sí mismo una firma producida cuando era válido, y el texto del reporte lo hace explícito, dejando la lectura jurídica al operador. Los errores ---flujo ilegible, PEM imparseable, CRL no confiable o vencida--- no son resultados: abortan la verificación.\par}

\begin{lstlisting}[language=Rust,caption={Veredicto de la verificación integral sobre los cuatro componentes (\texttt{crates/application/src/verification.rs}).},label={lst:impl-verdict}]
    let components = [&integrity, &signature, &certificate, &timestamp];
    let verdict = if components
        .iter()
        .any(|component| component.status == ComponentStatus::Failed)
    {
        Verdict::NotValid
    } else {
        Verdict::Valid
    };

    Ok(VerificationReport {
        document_digest_hex: digest.to_hex(),
        integrity,
        signature,
        certificate,
        timestamp,
        verdict,
    })
}
\end{lstlisting}

\subsection{S7. Hash de contraseñas con Argon2id}\label{sec:impl-s7-argon2}

El puerto \texttt{Password\allowbreak Hasher} (\rutarepo{crates/domain/src/crypto/password.rs}) declara el contrato del almacenamiento de contraseñas: cada hash usa una sal aleatoria fresca y produce una cadena en formato PHC, y una cadena almacenada ilegible es un error (\texttt{Malformed\allowbreak Password\allowbreak Hash}) que jamás se confunde con un intento de acceso fallido. El adaptador \texttt{Argon2id\allowbreak Hasher} (\rutarepo{crates/infrastructure/src/password.rs}) lo implementa con el crate \texttt{argon2}~0.5 y sales de \texttt{getrandom}, con los parámetros calibrados que fijó el diseño en \cref{sec:autenticacion} declarados como constantes (Listado~\ref{lst:impl-argon2}): 262\,144~KiB de memoria (256~MiB), dos iteraciones, un carril, sal de 16 bytes y salida de 32 bytes. El PHC resultante inicia con \texttt{\$argon2id\$v=19\$} y contiene \texttt{m=262144,t=2,p=1}, propiedad fijada por prueba. En la demostración de cierre, el hardware de referencia promedió 529.4 ms sobre cinco corridas, dentro de la banda objetivo inclusiva de 500 a 1\,000 milisegundos, y emite un veredicto \texttt{BelowBand}, \texttt{InBand} o \texttt{AboveBand}. Si el hardware de despliegue difiere, la medición debe repetirse.

\begin{lstlisting}[language=Rust,caption={Parámetros de Argon2id declarados como constantes (\texttt{crates/infrastructure/src/password.rs}).},label={lst:impl-argon2}]
/// Memory cost in KiB (256 MiB), calibrated on the reference host.
const MEMORY_KIB: u32 = 262144;

/// Number of passes over the memory.
const ITERATIONS: u32 = 2;

/// Degree of parallelism (lanes).
const PARALLELISM: u32 = 1;

/// Length in bytes of the derived hash.
const OUTPUT_LEN: usize = 32;

/// Length in bytes of the random salt drawn for every hash.
const SALT_LEN: usize = 16;
\end{lstlisting}

\subsection{S8. Segundo factor TOTP y códigos de recuperación}\label{sec:impl-s8-totp}

El puerto \texttt{Totp\allowbreak Provider} (\rutarepo{crates/domain/src/crypto/totp.rs}) declara el enrolamiento y la verificación del segundo factor; su adaptador \texttt{Totp\allowbreak Rs\allowbreak Provider} (\rutarepo{crates/infrastructure/src/totp.rs}) usa el crate \texttt{totp-rs}~5 con las features \texttt{otpauth} y \texttt{zeroize}, configurada con los parámetros del diseño (\cref{sec:autenticacion}): HMAC-SHA1, códigos de seis dígitos, paso de 30 segundos, desfase aceptado de un paso a cada lado (ventana efectiva de unos 90 segundos) y secreto de 20 bytes ---igual a la salida de HMAC-SHA1, como recomienda RFC 4226 \cite{mraihi2005}--- con un mínimo aceptado de 16 bytes \cite{mraihi2011}. El URI de aprovisionamiento \texttt{otpauth://} embebe el secreto codificado en base32 sin relleno \cite{josefsson2006}; por ello, el value object \texttt{Totp\allowbreak Enrollment} guarda sus tres campos ---secreto crudo, secreto en base32 y URI--- en buffers \texttt{Zeroizing}. La comparación de códigos ya es de tiempo constante dentro de \texttt{totp-rs}. El puerto aislado sigue siendo deliberadamente sin estado; en la aplicación HTTP, \texttt{Redis\allowbreak Session\allowbreak Store} materializa la decisión de \rutarepo{docs/adr/0008-totp-single-use-enforcement.md}: reclama atómicamente una huella del código aceptado con TTL de 90 segundos, por lo que una reproducción dentro de la ventana se rechaza.

{\emergencystretch=2em El enrolamiento completo lo orquesta el caso de uso \texttt{EnrollTotp}, que además emite los ocho códigos de recuperación de un solo uso previstos por el diseño. El generador \texttt{Random\allowbreak Recovery\allowbreak Code\allowbreak Generator} produce códigos de diez símbolos en la forma \texttt{XXXXX-XXXXX} sobre un alfabeto de 32 símbolos que excluye los caracteres confundibles I, L, O y U; cada símbolo aporta cinco bits extraídos sin sesgo de módulo, para un total de 50 bits de entropía por código. Los códigos en claro se muestran una única vez (viajan como \texttt{Zeroizing}); lo que se almacena es el value object \texttt{Recovery\allowbreak Code\allowbreak Set}: ocho ranuras con hashes PHC producidos por el mismo puerto \texttt{Password\allowbreak Hasher} de S7, cuya operación \texttt{consume} verifica un código contra las ranuras no usadas y consume la que coincide.\par}

\subsection{S9. Cadena de auditoría}\label{sec:impl-s9-auditoria}

La bitácora de auditoría es lógica de dominio pura (\rutarepo{crates/domain/src/audit/}). Cada evento registra secuencia, timestamp, actor, acción y recurso, y la cadena se define por la regla
\[
  \mathit{chain}(n) = \text{SHA-256}\bigl(\mathit{chain}(n-1) \,\Vert\, \text{bytes canónicos}(n)\bigr),
\]
con un valor génesis fijo de 32 bytes cero para la primera entrada. La codificación canónica serializa la secuencia como 8 bytes big-endian y cada campo variable con un prefijo de longitud de 4 bytes big-endian, de modo que dos eventos distintos jamás producen los mismos bytes; el timestamp se normaliza siempre a UTC y su forma RFC 3339 es la única representación textual que cubre la cadena. El Listado~\ref{lst:impl-chain} muestra el cálculo del eslabón. El hashing pasa por el puerto \texttt{Document\allowbreak Hasher}, así que ningún backend criptográfico entra al crate de dominio.

\begin{lstlisting}[language=Rust,caption={Regla de encadenamiento de la bitácora (\texttt{crates/domain/src/audit/chain.rs}).},label={lst:impl-chain}]
pub const GENESIS_PREVIOUS: Sha256Digest = Sha256Digest::from_array([0u8; SHA256_LEN]);

/// Computes the chain value of an event given the previous chain value.
///
/// # Errors
///
/// Propagates the canonical-encoding failures of
/// [`AuditEvent::canonical_bytes`].
pub fn chain_digest(
    hasher: &dyn DocumentHasher,
    previous: &Sha256Digest,
    event: &AuditEvent,
) -> Result<Sha256Digest, DomainError> {
    let canonical = event.canonical_bytes()?;
    let mut linked = Vec::with_capacity(SHA256_LEN + canonical.len());
    linked.extend_from_slice(previous.as_bytes());
    linked.extend_from_slice(&canonical);
    Ok(hasher.hash_bytes(&linked))
}
\end{lstlisting}

La verificación recomputa cada eslabón desde el génesis y señala el primer índice inconsistente. Detecta alteraciones de campos, inserciones y cambios de orden dentro de la cadena. No detecta por sí sola el truncamiento de la cola: un prefijo válido sigue verificando desde el mismo génesis. Se requiere un anclaje independiente de la cabeza para cubrir esa limitación, que permanece abierta en \rutarepo{docs/adr/0007-audit-chain-anchoring.md}.

El servidor usa \texttt{Postgres\allowbreak Audit\allowbreak Log} y comparte la cadena con las mutaciones de usuarios, expedientes, asignaciones y documentos. El adaptador transaccional calcula secuencia y hash bajo un bloqueo común; el cambio durable y su evento se confirman juntos. El timestamp se conserva como texto canónico, incluida su precisión de nanosegundos, para evitar que una conversión de almacenamiento altere los hashes. Los detalles de autorización, concurrencia y privilegios se presentan en \cref{sec:impl-http}.

\texttt{File\allowbreak Audit\allowbreak Log} permanece disponible para los comandos offline: guarda JSON Lines y coordina lectores y escritores con bloqueos sobre un archivo hermano. \texttt{In\allowbreak Memory\allowbreak Audit\allowbreak Log} sirve a pruebas y procesos efímeros. Los tres adaptadores conservan la regla del dominio; el puerto \texttt{Clock}, implementado por \texttt{SystemClock}, permite controlar el instante en pruebas. Los archivos de un origen importado quedan protegidos contra nuevas escrituras por los marcadores del procedimiento de corte.

\subsection{S10. Paquete de evidencia}\label{sec:impl-s10-evidencia}

El caso de uso \texttt{Export\allowbreak Evidence\allowbreak Package} (\rutarepo{crates/application/src/evidence.rs}) construye el paquete autocontenido que el diseño prometió en \cref{sec:sellado-tiempo}: un archivo ZIP con el documento, su firma separada (\texttt{<doc>.sig}), su sello de tiempo (\texttt{<doc>.tsr}), el certificado del firmante (\texttt{certificado.pem}), la raíz emisora (\texttt{ca.pem}), la CRL incluida como evidencia (\texttt{crl.pem}), opcionalmente la cadena de la autoridad de sellado (\texttt{tsa-chain.pem}) y un archivo \texttt{INSTRUCCIONES.md} generado desde una plantilla versionada, que documenta la verificación independiente completa en cuatro pasos con comandos literales de \texttt{openssl} ---incluida la versión local de la herramienta con que se redactaron---. La regla del ancla del sello es explícita: si \texttt{tsa-chain.pem} viaja en el paquete, la comprobación del token se ancla en esa cadena; en caso contrario, en \texttt{ca.pem}, pues la autoridad interna de sellado encadena a la misma raíz.

Los cuatro pasos de las instrucciones reproducen, con herramientas ajenas al prototipo, la verificación integral de S6:

\begin{itemize}
  \item \textbf{Integridad}: \texttt{openssl dgst -sha256} sobre el documento, comparado carácter por carácter con el resumen registrado en las propias instrucciones.
  \item \textbf{Firma}: extracción de la clave pública del certificado del firmante (\texttt{openssl x509 -pubkey}) y verificación de la firma separada con \texttt{openssl dgst -sha256 -verify}.
  \item \textbf{Certificado}: concatenación de \texttt{ca.pem} y \texttt{crl.pem}, seguida de \texttt{openssl verify} con la opción \texttt{-crl\_check}. Se identifican el instante de evaluación y la CRL suministrada. En el flujo HTTP se conserva la CRL capturada al sellar; exportar no la actualiza ni reemite la firma o el sello. El resultado técnico debe interpretarse según el material y el instante evaluados.
  \item \textbf{Sello de tiempo}: \texttt{openssl ts -verify} contra el ancla resuelta por la regla anterior, con la inspección del token disponible vía \texttt{openssl ts -reply -in <sello> -text}.
\end{itemize}

{\emergencystretch=2em El puerto \texttt{Archive\allowbreak Writer} lo implementa \texttt{Stored\allowbreak Zip\allowbreak Writer} (\rutarepo{crates/infrastructure/src/archive.rs}), una implementación propia del subconjunto mínimo del formato ZIP \cite{pkware2022}: entradas almacenadas sin compresión (método 0), un encabezado local por entrada, directorio central y registro de fin de directorio, con CRC-32 calculado bit a bit (polinomio reflejado \texttt{0xEDB88320}, verificado contra el valor de comprobación estándar \texttt{0xCBF43926}). La salida es determinista ---orden de entradas del llamador y timestamps fijos a la época DOS---, de modo que archivar dos veces el mismo contenido produce bytes idénticos. La alternativa de depender del crate \texttt{zip} se descartó en \rutarepo{docs/adr/0006-evidence-package-format.md}: el requisito de compilador fue una restricción del corte inicial, y su configuración por defecto incorporaba backends de compresión y criptografía innecesarios para el formato elegido, mientras que aquí la compresión no aporta nada ---se prioriza conservar bytes y evitar dependencias de compresión en el paquete de evidencia---. La seguridad de extracción se garantiza en el dominio: el value object \texttt{Archive\allowbreak Entry} restringe los nombres a un subconjunto ASCII conservador sin separadores de ruta, así que un archivo extraído jamás escribe fuera de su directorio destino.\par}

\section{La interfaz de línea de comandos}\label{sec:impl-cli}

El primer adaptador driving entregado por el prototipo es la interfaz de línea de comandos \texttt{despacho-cli}, construida con \texttt{clap}~4 en modo derive (\rutarepo{crates/bin/src/cli.rs}). El diseño anticipó un adaptador de línea de comandos para tareas administrativas junto a los adaptadores REST y WebSocket (\cref{sec:vista-componentes}); el orden de construcción privilegió la CLI porque permitió ejercitar primero los servicios criptográficos de extremo a extremo y después reutilizarlos desde HTTP. La bandera global \texttt{-{}-json} conmuta las operaciones administrativas a un objeto JSON legible por máquina. La \cref{tab:impl-cli-catalogo} presenta el catálogo completo de subcomandos con sus banderas reales, y la \cref{tab:impl-cli-casos} el cableado de cada grupo hacia los casos de uso de \texttt{application} y los adaptadores concretos que la raíz de composición inyecta.

\begingroup
\footnotesize
\begin{longtable}{@{}L{7.2cm}L{\dimexpr\linewidth-7.2cm-2\tabcolsep\relax}@{}}
  \caption{Catálogo de subcomandos de \texttt{despacho-cli}. La bandera global \texttt{-{}-json} aplica a todos.}
  \label{tab:impl-cli-catalogo} \\
    \toprule
    \textbf{Comando y banderas} & \textbf{Efecto} \\
    \midrule
    \endfirsthead
    \toprule
    \textbf{Comando y banderas} & \textbf{Efecto} \\
    \midrule
    \endhead

    \texttt{serve -{}-signer-cert <CERT> -{}-signer-key <KEY> [-{}-bind <ADDR>] [-{}-data-dir <DIR>]} & Inicia la API por expediente con un rol PostgreSQL restringido, sesiones Redis y TSA local. \\
    \midrule
    \texttt{database migrate -{}-runtime-role <ROLE>} & Aplica esquema y permisos con una conexión administrativa; el rol debe existir. \\
    \texttt{database import [-{}-data-dir <DIR>] -{}-mapping <PATH> [-{}-apply]} & Inspecciona origen y destino; solo \texttt{-{}-apply} confirma la importación y el corte. \\
    \midrule
    \texttt{crypto hash <FILE>} & Resumen SHA-256 del archivo. \\
    \midrule
    \texttt{vault encrypt <FILE> -{}-doc-id <UUID> [-{}-version <N>]} & Cifra con envelope encryption; escribe \texttt{<FILE>.enc}. \\
    \texttt{vault decrypt <PACKAGE> -{}-doc-id <UUID> [-{}-version <N>] [-{}-out <PATH>]} & Descifra un archivo \texttt{DVLT1}; sin \texttt{-{}-out} escribe el claro a la salida estándar. \\
    \texttt{vault rotate-kek -{}-file <PATH>} & Re-envuelve la DEK bajo la KEK de \texttt{NEW\_KEK\_BASE64}. \\
    \midrule
    \texttt{pki [-{}-scripts-dir <DIR>] init-ca} & Crea la CA raíz interna (idempotente). \\
    \texttt{pki issue -{}-cn <CN>} & Emite un certificado de entidad final. \\
    \texttt{pki revoke -{}-serial <HEX>} & Revoca por número de serie. \\
    \texttt{pki gen-crl} & Genera la CRL en PEM. \\
    \texttt{pki show <CERT>} & Inspecciona sujeto, emisor, serie y vigencia. \\
    \midrule
    \texttt{sign <FILE> -{}-cert <PATH> -{}-key <PATH>} & Firma el digest del archivo; escribe \texttt{<FILE>.sig}. \\
    \texttt{timestamp <INPUT> [-{}-mock] [-{}-out <PATH>] [-{}-pki-dir <DIR>]} & Solicita un sello RFC 3161; escribe \texttt{<INPUT>.tsr}. Con \texttt{-{}-mock} usa la TSA local; sin ella, el proveedor remoto. \\
    \texttt{verify <FILE> -{}-sig <PATH> -{}-cert <PATH> -{}-ca <PATH> [-{}-tsr] [-{}-crl] [-{}-tsa-cert] [-{}-at <RFC3339>]} & Verificación integral de cuatro componentes; \texttt{-{}-at} evalúa el certificado a un instante dado. \\
    \texttt{package export <FILE> -{}-sig -{}-tsr -{}-cert -{}-ca -{}-crl [-{}-tsa-chain] -{}-out <PATH>} & Exporta el paquete de evidencia, previa verificación de los artefactos. \\
    \midrule
    \texttt{auth calibrate} & Mide el costo real de Argon2id contra la banda de 500--1\,000 ms. \\
    \texttt{auth hash-password} & Hash PHC de una contraseña leída de la entrada estándar. \\
    \texttt{auth verify-password -{}-hash <PHC>} & Verifica una contraseña contra un PHC almacenado. \\
    \texttt{auth totp enroll -{}-user <EMAIL> -{}-secret-out <PATH>} & Enrola el segundo factor y emite los códigos de recuperación. \\
    \texttt{auth totp verify -{}-code <CODE> -{}-secret-file <PATH>} & Verifica un código TOTP. \\
    \midrule
    \texttt{audit append -{}-action <STRING> -{}-resource <STRING>} & Anexa un evento a la bitácora offline, si el origen no está bloqueado por migración. \\
    \texttt{audit verify-chain} & Verifica la cadena del archivo offline desde el génesis. \\
    \texttt{audit show} & Lista los eventos del archivo offline en orden de inserción. \\
    \bottomrule
\end{longtable}
\endgroup

\begin{table}[htbp]
  \centering
  \caption{Correspondencia entre comandos, orquestación y adaptadores compuestos por el binario.}
  \label{tab:impl-cli-casos}
  \begin{tabularx}{\linewidth}{@{}L{3.4cm} L{5.2cm} >{\raggedright\arraybackslash}X@{}}
    \toprule
    \textbf{Comando} & \textbf{Caso(s) de uso} & \textbf{Adaptadores inyectados} \\
    \midrule
    \texttt{serve} & \texttt{Identity\allowbreak Service}, \texttt{CaseService}, \texttt{CaseDocument\allowbreak Service} & PostgreSQL transaccional, Redis, Argon2id, TOTP, cifrado, firma, TSA local y ZIP \\
    \texttt{database *} & Preparación administrativa y validación de registros históricos & \texttt{initialize\_database}, \texttt{Legacy\allowbreak Import} y puertos de validación criptográfica \\
    \texttt{crypto hash} & \texttt{Hash\allowbreak Document} & \texttt{Ring\allowbreak Sha256\allowbreak Hasher} \\
    \texttt{vault *} & \texttt{Encrypt\allowbreak Document}, \texttt{Decrypt\allowbreak Document}, \texttt{RotateKek} & \texttt{Ring\allowbreak Aes\allowbreak Gcm\allowbreak Cipher}, \texttt{Envelope\allowbreak Key\allowbreak Manager} \\
    \texttt{pki *} & \texttt{Initialize\allowbreak Ca}, \texttt{Issue\allowbreak Certificate}, \texttt{Revoke\allowbreak Certificate}, \texttt{GenerateCrl}, \texttt{Inspect\allowbreak Certificate} & \texttt{Openssl\allowbreak Ca\allowbreak Adapter}, \texttt{X509\allowbreak Chain\allowbreak Validator} \\
    \texttt{sign} & \texttt{Sign\allowbreak Document} + \texttt{Verify\allowbreak Signature} & \texttt{Rsa\allowbreak Pkcs1\allowbreak Signer}, \texttt{Rsa\allowbreak Pkcs1\allowbreak Verifier} \\
    \texttt{timestamp} & \texttt{Timestamp\allowbreak Document} & \texttt{Local\allowbreak Openssl\allowbreak Tsa} o \texttt{Cincel\allowbreak Tsa\allowbreak Adapter} \\
    \texttt{verify} & \texttt{Verify\allowbreak Document} & \texttt{Ring\allowbreak Sha256\allowbreak Hasher}, \texttt{Rsa\allowbreak Pkcs1\allowbreak Verifier}, \texttt{X509\allowbreak Chain\allowbreak Validator}, \texttt{Rfc3161\allowbreak Verifier} \\
    \texttt{package export} & \texttt{Verify\allowbreak Document} + \texttt{Export\allowbreak Evidence\allowbreak Package} & los cuatro anteriores más \texttt{Stored\allowbreak Zip\allowbreak Writer} \\
    \texttt{auth calibrate}, \texttt{hash-password}, \texttt{verify-password} & \texttt{CalibratePassword\allowbreak Hashing}, \texttt{Hash\allowbreak Password}, \texttt{Verify\allowbreak Password} & \texttt{Argon2id\allowbreak Hasher} \\
    \texttt{auth totp *} & \texttt{EnrollTotp}, \texttt{VerifyTotp} & \texttt{Totp\allowbreak Rs\allowbreak Provider}, \texttt{Random\allowbreak Recovery\allowbreak Code\allowbreak Generator}, \texttt{Argon2id\allowbreak Hasher} \\
    \texttt{audit *} & \texttt{Append\allowbreak Audit\allowbreak Event}, \texttt{Verify\allowbreak Audit\allowbreak Chain}, \texttt{Show\allowbreak Audit\allowbreak Log} & \texttt{File\allowbreak Audit\allowbreak Log}, \texttt{Ring\allowbreak Sha256\allowbreak Hasher}, \texttt{SystemClock} \\
    \bottomrule
  \end{tabularx}
\end{table}

Con \texttt{-{}-json}, las operaciones finitas emiten resultados estructurados en la salida estándar ---por ejemplo, \texttt{crypto hash} emite \texttt{\{algorithm, digest, file\}} y \texttt{verify} emite el reporte completo con los cuatro componentes y el veredicto---, mientras los diagnósticos van a la salida de error. Se exceptúan el proceso persistente \texttt{serve} y \texttt{vault decrypt} sin archivo de salida, que conserva el texto claro crudo. Los códigos de salida están pensados para que los scripts bifurquen directamente: el proceso termina con código distinto de cero cuando \texttt{verify} produce un veredicto no válido, cuando \texttt{auth verify-password} reporta una contraseña que no coincide, cuando \texttt{auth totp verify} rechaza el código, cuando \texttt{audit verify-chain} encuentra la cadena rota ---imprimiendo antes el JSON de fallo con el índice del primer eslabón roto--- y cuando \texttt{package export} rehúsa exportar. Los subcomandos de lectura de la bitácora (\texttt{verify-chain}, \texttt{show}) rehúsan operar sobre un archivo inexistente: reportar una ruta equivocada o un log borrado como una cadena vacía válida frustraría el propósito de la verificación.

La raíz de composición (\rutarepo{crates/bin/src/main.rs}) es el único lugar del sistema donde los adaptadores concretos se instancian e inyectan en los casos de uso, y el único crate que usa \texttt{anyhow}, tal como estableció el patrón de inyección de dependencias por composición del diseño (\cref{sec:inyeccion-dependencias}). Ahí ---y no en los adaptadores, que quedan libres de acceso al entorno--- se leen los secretos: la KEK de \texttt{KEK\_BASE64} a un buffer \texttt{Zeroizing}, la clave privada del firmante desde su archivo PEM, la credencial del proveedor de sellado. La raíz concentra también salvaguardas operativas que no pertenecen a ningún servicio individual: \texttt{sign} verifica la firma recién producida contra el certificado \emph{antes} de escribir \texttt{<FILE>.sig}, de modo que una clave que no corresponde al certificado falla ahí y no produce una firma inverificable; \texttt{timestamp} relee el token recibido y comprueba que cubra el digest calculado antes de escribir el archivo \texttt{.tsr}; \texttt{vault rotate-kek} escribe a un archivo temporal y renombra encima, de modo que un fallo a la mitad deja el original intacto; y \texttt{package export} somete los artefactos al mismo caso de uso \texttt{Verify\allowbreak Document} de la verificación integral ---CRL y token incluidos--- y rehúsa exportar si el veredicto no es válido, nombrando cada componente fallido, porque un paquete que fallaría sus propias instrucciones de verificación no debe existir.

La integración de todos los servicios se demuestra con el guion \rutarepo{scripts/demo.sh}, que compila el binario y recorre el ciclo completo de evidencia en un directorio temporal: creación de la autoridad interna, del firmante y de la autoridad de sellado; hash del documento de muestra; cifrado en reposo con corrupción deliberada de un byte y observación del rechazo; firma y obtención del sello de tiempo; verificación integral con los cuatro componentes válidos; alteración de una copia del documento con el consecuente fallo del componente de firma; exportación del paquete de evidencia y su verificación con solo \texttt{openssl} y \texttt{unzip}; revocación del certificado y su detección; calibración de autenticación y códigos de un solo uso; y bitácora encadenada con anexado, verificación, manipulación y nueva verificación. La transcripción íntegra de una ejecución real, capturada el 12 de julio de 2026 sobre OpenSSL 3.5.7, está versionada en \rutarepo{docs/demo-transcript.md}. Una segunda demostración, \rutarepo{scripts/api-demo.sh}, recorre el mismo núcleo a través de HTTP y se describe en la sección siguiente. El entorno, los comandos y el resumen de ambos recorridos quedaron asentados en \rutarepo{docs/verification-report.md}.

\section{Aplicación HTTP autenticada y persistencia transaccional}\label{sec:impl-http}

La capa HTTP delega cada operación en un puerto de aplicación. \texttt{Identity\allowbreak Workflow} se implementa en \texttt{Identity\allowbreak Service} para autenticación; \texttt{Case\allowbreak Workflow}, en \texttt{CaseService} para expedientes; y \texttt{Case\allowbreak Document\allowbreak Workflow}, en \texttt{Case\allowbreak Document\allowbreak Service} para documentos. \texttt{Document\allowbreak Processor} concentra la preparación criptográfica compartida con el flujo offline; el servidor no inyecta el repositorio documental de archivos. La raíz de composición selecciona \texttt{Postgres\allowbreak Case\allowbreak Document\allowbreak Store}, la auditoría PostgreSQL y \texttt{Local\allowbreak Openssl\allowbreak Tsa}. La ejecución HTTP no consulta Cincel ni sustituye silenciosamente la TSA por un proveedor remoto.

\subsection{Identidad y sesiones vigentes}\label{sec:impl-http-identidad}

El alta inicial solo funciona mientras no existen usuarios. PostgreSQL serializa esa condición y confirma el primer Owner junto con su evento de auditoría. El correo se normaliza y la contraseña se almacena como PHC Argon2id; el enrolamiento entrega una sola vez el secreto TOTP y ocho códigos de recuperación. La tabla \texttt{users} conserva UUID, correo único, rol, estado activo, secreto TOTP cifrado, hashes de recuperación y revisión optimista. \texttt{Aes\allowbreak Gcm\allowbreak Secret\allowbreak Protector} liga el secreto al UUID del usuario mediante datos autenticados bajo la KEK del entorno.

Login verifica la contraseña y crea un desafío Redis con duración de cinco minutos. Cada intento de segundo factor toma ese desafío con \texttt{GETDEL}, operación atómica que impide que dos factores emitan sesiones a partir del mismo desafío; un rechazo también lo consume. TOTP utiliza además un reclamo atómico contra repetición. El consumo de recuperación actualiza el conjunto y su evento dentro de PostgreSQL. Antes de emitir una sesión se vuelve a comprobar que la cuenta esté activa. Las sesiones usan tokens opacos de 256 bits, con TTL de 24 horas; Redis conserva su índice SHA-256, no el token en claro. Cinco contraseñas rechazadas bloquean durante 15 minutos la clave normalizada del correo. Incremento y expiración del contador son atómicos.

\begin{table}[H]
  \centering
  \caption{Contrato HTTP de identidad.}
  \label{tab:impl-http-identidad}
  \begin{tabularx}{\linewidth}{@{}L{6.2cm} X@{}}
    \toprule
    \textbf{Ruta} & \textbf{Responsabilidad} \\
    \midrule
    \texttt{POST /api/v1/auth/bootstrap} & Crea el primer Owner y entrega su enrolamiento; \texttt{201}. \\
    \texttt{POST /api/v1/auth/login} & Verifica contraseña y entrega desafío MFA. \\
    \texttt{POST /api/v1/auth/mfa/totp} & Consume desafío, verifica TOTP y crea sesión. \\
    \texttt{POST /api/v1/auth/mfa/recovery} & Consume desafío y recuperación; crea sesión. \\
    \texttt{GET /api/v1/auth/me} & Devuelve identidad y rol vigentes. \\
    \texttt{POST /api/v1/auth/logout} & Revoca la sesión; \texttt{204}. \\
    \texttt{POST /api/v1/users} & Owner crea y enrola otra cuenta; \texttt{201}. \\
    \bottomrule
  \end{tabularx}
\end{table}

Cada operación protegida resuelve la sesión en Redis y recarga estado y rol desde PostgreSQL. La creación de usuarios autentica dentro del caso de uso y revalida al Owner al confirmar. PostgreSQL y Redis no comparten una transacción distribuida: si falla la auditoría tras crear un desafío o sesión, la aplicación intenta retirar la credencial y devuelve error sin token. Si también falla la limpieza, la clave puede permanecer hasta su expiración. Logout revoca primero y no restaura la sesión si falla el evento posterior. La pérdida de la respuesta de enrolamiento tampoco dispone todavía de un procedimiento completo de recuperación del ciclo de alta.

\subsection{Expedientes, pertenencia y aislamiento}\label{sec:impl-http-expedientes}

El dominio identifica cada expediente mediante \texttt{CaseId} y valida título y referencia con \texttt{Case\allowbreak Metadata}: ambos son obligatorios, se recortan espacios exteriores, se rechazan controles y se limitan a 200 y 100 caracteres, respectivamente. La referencia es una etiqueta libre, no un identificador único. PostgreSQL conserva metadatos, creador y asignaciones de usuarios. Estas asignaciones regulan acceso; todavía no representan partes procesales, audiencias ni plazos.

\begin{table}[H]
  \centering
  \caption{Contrato HTTP de expedientes y asignaciones. Los parámetros \texttt{<case>} y \texttt{<user>} son UUID.}
  \label{tab:impl-http-expedientes}
  \begin{tabularx}{\linewidth}{@{}L{6.2cm} X@{}}
    \toprule
    \textbf{Ruta} & \textbf{Responsabilidad} \\
    \midrule
    \texttt{POST /api/v1/cases} & Owner o Litigante crean y quedan asignados; \texttt{201}. \\
    \texttt{GET /api/v1/cases} & Lista únicamente metadatos visibles, con \texttt{limit} y \texttt{offset}. \\
    \texttt{GET /api/v1/cases/<case>} & Devuelve metadatos si el actor es Owner o miembro vigente. \\
    \texttt{PUT /api/v1/cases/\allowbreak <case>/members/<user>} & Owner asigna un usuario activo; operación idempotente, \texttt{204}. \\
    \texttt{DELETE /api/v1/cases/\allowbreak <case>/members/<user>} & Owner retira pertenencia; operación idempotente, \texttt{204}. \\
    \bottomrule
  \end{tabularx}
\end{table}

Owner ve todos los expedientes y administra asignaciones. Litigante puede crear un expediente, queda asignado automáticamente y pierde acceso si se retira esa asignación. Paralegal y Cliente consultan solamente metadatos asignados. La consulta filtra pertenencia antes de paginar; el límite predeterminado es 50, admite de 1 a 100 y el desplazamiento es no negativo. El orden es por UUID, sin promesa de una instantánea común entre páginas. Un expediente ajeno responde igual que uno inexistente: \texttt{404 case\_not\_found}.

\subsection{Operaciones documentales por expediente}\label{sec:impl-http-documentos}

La carga recibe hasta 16~MiB binarios y un nombre seguro en \texttt{X-Document-Name}. El actor deriva de \texttt{Authorization: Bearer}; \texttt{X-Actor} no forma parte del contrato. El servicio crea UUID y versión uno, calcula SHA-256 y prepara el paquete cifrado \texttt{DVLT1}. El AAD conserva el par identidad documental--versión; la asociación con el expediente es una restricción inmutable de PostgreSQL, no un campo agregado retroactivamente al contexto criptográfico.

\begin{table}[H]
  \centering
  \caption{Contrato documental por expediente. Todas las rutas de las cuatro primeras filas usan el prefijo \texttt{/api/v1/cases/<case>/documents}.}
  \label{tab:impl-http-rutas}
  \begin{tabularx}{\linewidth}{@{}L{5.2cm} X@{}}
    \toprule
    \textbf{Método y sufijo} & \textbf{Responsabilidad} \\
    \midrule
    \texttt{POST} al prefijo & Carga cifrada y asociada al expediente; \texttt{201}. \\
    \texttt{POST /<document>/seal} & Firma y sello local; confirma evidencia una sola vez. \\
    \texttt{POST /<document>/verify} & Reporte de integridad, firma, certificado/CRL y sello. \\
    \texttt{GET /<document>/evidence} & ZIP de evidencia y cabecera \texttt{X-Document-Digest}. \\
    \midrule
    \texttt{GET /api/v1/audit/verify} & Owner comprueba la cadena PostgreSQL completa. \\
    \texttt{GET /healthz} & Diagnóstico público del proceso. \\
    \bottomrule
  \end{tabularx}
\end{table}

Owner accede a cualquier expediente. Litigante y Paralegal necesitan pertenencia vigente para cargar, verificar o exportar; solamente Owner y Litigante pueden sellar. Cliente conserva denegadas todas las acciones documentales incluso cuando está asignado. El rol se comprueba antes de revelar el recurso. Un documento fuera del expediente indicado responde \texttt{404 document\_not\_found}, incluso para Owner o un usuario perteneciente a ambos expedientes. Las antiguas rutas globales bajo \texttt{/api/v1/documents} responden \texttt{404} y no tienen alias.

Carga y sellado devuelven un objeto con \texttt{case\_id}, \texttt{id}, \texttt{version}, \texttt{name}, \texttt{digest} y \texttt{sealed}. No existen todavía listado, detalle, búsqueda ni historial de versiones documentales. Sellar requiere descifrar y comprobar el digest, producir y verificar la firma, solicitar el token y verificarlo bajo la cadena TSA capturada. Una segunda solicitud sobre un documento sellado produce \texttt{409 document\_already\_sealed}. Exportar comprueba el contenido descifrado y empaqueta la evidencia capturada sin actualizarla; la consulta de verificación proporciona por separado el reporte criptográfico completo.

\subsection{Confirmación transaccional y continuidad de evidencia}\label{sec:impl-http-transacciones}

El esquema durable se organiza en tres migraciones:
\begin{itemize}
  \item \rutarepo{migrations/0001_identity.sql}: usuarios y credenciales protegidas.
  \item \rutarepo{migrations/0002_cases.sql}: expedientes y asignaciones.
  \item \rutarepo{migrations/0003_case_documents_audit.sql}: documentos, auditoría y recibos de importación.
\end{itemize}

El servicio autentica antes de la preparación y de nuevo inmediatamente antes de confirmar, exigiendo la misma identidad. El adaptador \texttt{Case\allowbreak Document\allowbreak Store} revalida usuario activo, rol, pertenencia y coincidencia exacta entre expediente y documento dentro de la transacción. Las operaciones costosas ---AES, RSA, TSA y ZIP--- se realizan fuera de ella.

Las transacciones mutantes del backend toman un mismo bloqueo consultivo antes de bloquear las filas pertinentes. Este orden serializa las confirmaciones y la cabeza de auditoría. Si una retirada de pertenencia confirma primero, la operación documental posterior se rechaza aun usando la misma sesión; si la operación documental obtiene primero el bloqueo, la retirada espera. Una respuesta confirmada puede terminar de transmitirse después de la revocación. La comprobación de la sesión Redis y la transacción PostgreSQL siguen siendo pasos separados, por lo que no se afirma atomicidad distribuida.

Documento y evento se confirman juntos. La creación de expediente y su asignación inicial, los cambios de pertenencia y las mutaciones durables de identidad también conservan su evento en la misma transacción. Verificación y exportación vuelven a comprobar acceso y estado documental y registran su evento antes de entregar el resultado. UUID, expediente, versión, nombre, digest, vault y evidencia ya sellada quedan protegidos por permisos de columna y un trigger. Dos preparaciones de sellado pueden competir y solicitar tokens; únicamente una persiste, y la otra recibe conflicto sin sustituir la evidencia ganadora.

\subsection{Preparación administrativa e importación}\label{sec:impl-http-migracion}

El esquema se prepara explícitamente con \texttt{database migrate -{}-runtime-role}, usando una conexión administrativa y un rol operativo previamente creado. \texttt{serve} no ejecuta DDL: valida que su rol no pueda administrar o reescribir auditoría, documentos, recibos ni el mecanismo que protege evidencia. El rol operativo puede leer y anexar eventos, pero no actualizarlos, eliminarlos ni truncarlos. Los administradores y propietarios de PostgreSQL permanecen dentro de la base de confianza.

La migración del repositorio local requiere detener escritores, conservar un respaldo y proporcionar un mapa completo entre cada documento y un expediente existente. \texttt{database import} inspecciona archivos y destino sin escribir; \texttt{-{}-apply} confirma. La inspección valida identidad, AAD, digest, firma, sello y cadena de auditoría. La cadena y CRL del firmante se evalúan en el instante autenticado del sello; la confianza en la TSA mantiene la política actual del verificador OpenSSL, de modo que una TSA expirada al ejecutar puede impedir importar evidencia histórica. Un rechazo conserva el origen para resolución manual, sin reemitir firmas ni sellos.

La primera importación exige vacíos los destinos de documentos y auditoría. Conserva la historia como prefijo, inserta documentos, añade el evento de importación y guarda el recibo en una transacción. Antes de insertar instala barreras \texttt{.migration-pending} en ambos almacenes de archivos; tras confirmar instala los marcadores \texttt{.migrated}. Si falla el corte de archivos después del commit, el error identifica la importación como confirmada y las fuentes permanecen bloqueadas. Repetir las mismas fuentes y mapa reconcilia el recibo y completa los marcadores sin duplicar datos. Los ejecutables antiguos no reconocen estas barreras y deben permanecer detenidos.

\texttt{serve -{}-data-dir} identifica el origen preservado, no el destino de documentos nuevos. El arranque reconcilia marcadores, recibo, asociaciones, documentos y cadena completa; un recibo aislado no basta para aceptar una restauración parcial. El respaldo operativo debe incluir PostgreSQL completo y proteger KEK y material criptográfico por separado. El procedimiento de importación y recuperación se desarrolla en \cref{sec:cli-database} y en \rutarepo{docs/database-operations.md}.

\subsection{Límites de ejecución y verificación integrada}\label{sec:impl-http-limites}

Los adaptadores PostgreSQL y Redis son síncronos. HTTP ejecuta los casos de uso en trabajadores bloqueantes de Tokio con un presupuesto compartido: ocho peticiones admitidas y dos operaciones bloqueantes por defecto, configurables mediante \texttt{-{}-max-in-flight-requests} y \texttt{-{}-max-blocking-operations}. La cancelación del cliente no libera el permiso hasta que termina el trabajo subyacente ni revierte una mutación ya confirmada. Sin cupo se responde \texttt{503 server\_busy}; \texttt{/healthz} permanece fuera de admisión. Los cuerpos JSON de identidad y expedientes se limitan a 16~KiB y los documentos a 16~MiB. El consumo por hash Argon2id obliga a dimensionar la concurrencia con mediciones del entorno.

Los errores de credenciales o sesión usan \texttt{401}; permisos, \texttt{403}; ausencia o aislamiento, \texttt{404}; conflictos, \texttt{409}; entradas inválidas, \texttt{400} o \texttt{422}; límite de login, \texttt{429}; y fallos internos sin secretos, \texttt{500}. El contrato se documenta en \rutarepo{docs/http-api.md}; sus límites y transacciones, en \rutarepo{docs/adr/0015-backend-concurrency-and-invariants.md} y \rutarepo{docs/adr/0016-case-document-transactions.md}. Quedan pendientes TLS, límites de transporte, pooling, plazos operativos completos, pruebas de carga, apagado ordenado y anclaje externo de auditoría.

\rutarepo{scripts/api-demo.sh} crea servicios desechables, usuarios de los cuatro roles, expedientes y documentos. Comprueba autenticación, permisos, consultas filtradas, retirada de pertenencia con la misma sesión y ausencia de accesos documentales del Cliente. Recorre firma, sello, exportación y verificación independiente con OpenSSL; ensaya también importación y restauración con evidencia sellada. \rutarepo{scripts/test-backends.sh} prepara las bases aisladas y Redis para pruebas reales. Sus ejecuciones y resultados se distinguen de la descripción funcional en el Capítulo~\ref{cap:pruebas}.

\section{Resultados de la exploración del sellado de tiempo}\label{sec:impl-spike-tsa}

El análisis de riesgos identificó la integración con la API del PSC como el riesgo técnico de mayor exposición (RTC-01, \cref{tab:rtc-01}), y previó como contingencia un servicio sustituto que simulara las respuestas de la TSA para no bloquear el desarrollo. El registro en el entorno de Cincel pudo completarse, pero el riesgo se materializó durante el consumo: la API no mostró estabilidad suficiente para sostener una campaña repetible y el sandbox requiere pago, por lo que no es una dependencia determinista ni económicamente neutra para la suite o la demostración (\cref{sec:riesgo-psc-materializado}). La respuesta de implementación fue más fuerte que la contingencia prevista: en lugar de un simulador, se construyó una \emph{autoridad de sellado real} operada localmente con OpenSSL, que emite tokens RFC 3161 auténticos, criptográficamente completos y verificables con herramientas independientes.

La autoridad local se prepara con el script \rutarepo{pki/issue-tsa-cert.sh}: la CA interna emite un certificado de TSA RSA-3072 con vigencia de 365 días cuya extensión \texttt{extended\allowbreak Key\allowbreak Usage} es crítica y contiene únicamente \texttt{time\allowbreak Stamping}, como exige RFC 3161 \cite{adams2001}. La configuración \rutarepo{pki/tsa.cnf} fija los parámetros de operación: digests de solicitud y de firma SHA-256, exactitud declarada de un segundo (\texttt{accuracy = secs:1}), sin garantía de orden entre tokens, números de serie desde 1000, y una política de emisión identificada por el OID 1.3.6.1.4.1.13762.3 ---un arco privado del árbol de empresas de IANA, declarado explícitamente como política de demostración y no como política de producción registrada; las solicitudes que pidan cualquier otra política se rehúsan---. El adaptador \texttt{Local\allowbreak Openssl\allowbreak Tsa} ejercita esta autoridad vía el comando \texttt{timestamp -{}-mock}, y las pruebas de integración \rutarepo{crates/infrastructure/tests/local_tsa_rfc3161.rs} recorren el ciclo completo contra los scripts reales en directorios temporales, cotejando los tokens emitidos con la línea de comandos de \texttt{openssl}; ninguna de esas pruebas está marcada como ignorada. La demostración de extremo a extremo del prototipo emite sus sellos con esta autoridad, y la verificación independiente del paquete de evidencia (\texttt{openssl ts -verify -data documento -in documento.tsr -CAfile ca.pem}) los acepta.

El camino remoto hacia el PSC no se eliminó: permanece como una integración explícita del comando \texttt{timestamp} cuando no se pasa \texttt{-{}-mock}. En ese modo, el binario exige \texttt{CINCEL\_BASE\_URL} y \texttt{CINCEL\_API\_KEY} no vacías en el entorno ---los valores en blanco de la plantilla no configuran nada--- y construye el \texttt{Cincel\allowbreak Tsa\allowbreak Adapter} con un tiempo límite de 15 segundos por solicitud y una política de reintento de cinco sondeos cada dos segundos. Debe decirse con precisión qué estatus tiene ese adaptador: su contrato HTTP ---envío del digest por \texttt{POST}, estados \texttt{completed}, \texttt{processing} y \texttt{rejected}, sondeo del recurso diferido--- es un contrato asumido por el adaptador y cubierto por un stub local, no un contrato confirmado por una campaña externa. El stub (\rutarepo{crates/infrastructure/tests/cincel_stub.rs}) cubre el token inmediato, el procesamiento diferido, los rechazos con detalle, los estados HTTP de error, los cuerpos indecodificables, la autoridad inalcanzable y la promesa de que la credencial jamás se filtra a los mensajes de error. La prueba de humo contra el sandbox real, \texttt{the\_real\_sandbox\_\allowbreak issues\_a\_token}, permanece ignorada: la API no ofreció garantía operativa suficiente y el sandbox requiere pago con precios no transparentes, por lo que no forma parte de la evidencia de esta entrega.

El balance honesto de la exploración es entonces el siguiente. El acceso al proveedor no produjo una ejecución suficientemente estable y repetible para declarar confirmado su contrato HTTP ni para medir el criterio de éxito de OE-3; tampoco existe un precio transparente que permita autorizar una campaña. Por decisión del proyecto, la confirmación empírica del algoritmo de firma de la TSA comercial y la tasa externa quedan fuera del alcance actual. La limitación quedó \emph{acotada arquitectónicamente}: el flujo RFC 3161 completo ---solicitud, emisión, almacenamiento del token DER y verificación en dos etapas--- está demostrado de extremo a extremo con una autoridad real, el verificador \texttt{Rfc3161\allowbreak Verifier} es agnóstico de la autoridad emisora porque verifica cualquier token contra el ancla de confianza suministrada, y el intercambio de emisor sigue siendo una sustitución explícita detrás del puerto \texttt{Timestamp\allowbreak Service}, sin tocar el dominio ni la lógica de verificación. Esta portabilidad no equipara las garantías: un token de la TSA interna prueba el funcionamiento técnico bajo la CA del prototipo, pero no constituye una constancia independiente emitida por un PSC autorizado.

\section{Divergencias respecto al diseño y decisiones registradas}\label{sec:impl-divergencias}

Durante la construcción, las decisiones con consecuencias arquitectónicas se registraron como Architecture Decision Records versionados en el directorio \rutarepo{docs/adr/} del repositorio, en el formato de Nygard \cite{nygard2011} ---contexto, decisión, estatus y consecuencias---. Estos registros son distintos y complementarios de los ADR de diseño presentados en \cref{sec:adrs}: aquellos capturaron las decisiones de arquitectura previas a la construcción; estos capturan las decisiones tomadas al materializarla, y el código y la configuración los citan por su ruta. La \cref{tab:impl-adrs} los resume.

\begingroup
\footnotesize
\begin{longtable}{@{}L{6.2cm}L{\dimexpr\linewidth-6.2cm-2\tabcolsep\relax}@{}}
  \caption{Architecture Decision Records registrados en \texttt{docs/adr/} durante la implementación.}
  \label{tab:impl-adrs} \\
    \toprule
    \textbf{Archivo} & \textbf{Decisión} \\
    \midrule
    \endfirsthead
    \toprule
    \textbf{Archivo} & \textbf{Decisión} \\
    \midrule
    \endhead

    \rutarepo{0001-msrv-and-dependency-pinning.md} & Pin inicial de los crates criptográficos a la familia \texttt{der}~0.7 (\texttt{x509-cert}~0.2, \texttt{cms}~0.2, \texttt{x509-tsp}~0.1, \texttt{rsa}~0.9; \texttt{ring} $\geq$ 0.17.12). \\
    \rutarepo{0002-rsa-signing-crate-and-advisory.md} & Firmar con el crate \texttt{rsa}; aviso RUSTSEC-2023-0071 registrado bajo el modelo inicial de amenaza local y declarado en \texttt{deny.toml} y \texttt{audit.toml}. \\
    \rutarepo{0003-zeroize-secret-material-in-domain.md} & Admitir \texttt{zeroize} en el dominio y declarar todo campo portador de secretos como buffer \texttt{Zeroizing}. \\
    \rutarepo{0004-certificate-validation-implementation.md} & Validar X.509 re-codificando a DER la estructura to-be-signed parseada y verificando la firma con \texttt{rsa} y \texttt{sha2}. \\
    \rutarepo{0005-rfc3161-verification-strategy.md} & Verificar tokens RFC 3161 en dos etapas: decodificación e imprint nativos; firma CMS y cadena delegadas a \texttt{openssl ts -verify}. \\
    \rutarepo{0006-evidence-package-format.md} & Escribir el paquete de evidencia con un writer ZIP propio de entradas almacenadas, determinista y sin dependencia nueva. \\
    \rutarepo{0007-audit-chain-anchoring.md} & Mantener la regla de cadena y documentar que el truncamiento de cola no es detectable; el anclaje de la cabeza se delega a la capa de persistencia. \\
    \rutarepo{0008-totp-single-use-enforcement.md} & Mantener el puerto TOTP sin estado y exigir un solo uso en el llamador mediante persistencia por cuenta. \\
    \rutarepo{0009-local-timestamp-authority.md} & Excluir la campaña externa no reproducible y demostrar RFC 3161 con una TSA OpenSSL local, sin atribuirle equivalencia jurídica. \\
    \rutarepo{0010-http-adapter-foundation.md} & Introducir Axum como adaptador driving, conservar la dirección de dependencias y comenzar por el contrato de salud. \\
    \rutarepo{0011-local-document-workflow.md} & Orquestar el flujo documental detrás de un puerto, persistir solo contenido cifrado y seleccionar explícitamente la TSA local en la aplicación. \\
    \rutarepo{0012-revocable-sessions-and-rbac.md} & Usar tokens opacos revocables, usuarios PostgreSQL, estado efímero Redis y RBAC conservador sin depender de Cincel. \\
    \rutarepo{0013-raise-msrv-for-security-fixes.md} & Elevar MSRV a Rust 1.88 para incorporar correcciones de seguridad y hacer obligatoria su comprobación. \\
    \rutarepo{0014-case-membership-and-isolation.md} & Persistir expedientes y asignaciones, filtrar lecturas y mantener denegado el acceso documental del Cliente. \\
    \rutarepo{0015-backend-concurrency-and-invariants.md} & Consumir desafíos MFA atómicamente, autenticar altas en aplicación y limitar trabajo HTTP bloqueante. \\
    \rutarepo{0016-case-document-transactions.md} & Autorizar documentos por expediente, confirmar cambios y auditoría en PostgreSQL e importar archivos con reconciliación. \\
    \bottomrule
\end{longtable}
\endgroup

Las divergencias sustantivas entre lo diseñado en el Capítulo~\ref{cap:analisis} y lo construido son las siguientes; se listan de forma explícita porque un documento que las omitiera daría una imagen más limpia pero menos útil del proceso de ingeniería.

\begin{itemize}
  \item \textbf{Biblioteca de firma.} El diseño situó los servicios criptográficos, de manera general, sobre la biblioteca \texttt{ring} (\cref{sec:hexagonal-definicion}). En la implementación, \texttt{ring} quedó a cargo de SHA-256 y AES-256-GCM, pero la firma se materializó con el crate \texttt{rsa} ---el adaptador conserva el nombre \texttt{Rsa\allowbreak Pkcs1\allowbreak Signer} que el diseño ya anticipaba en \cref{sec:inyeccion-dependencias}--- porque interopera directamente con las claves y certificados PEM que produce la CA interna vía \texttt{x509-cert} y \texttt{sha2}. El costo fue aceptar el aviso RUSTSEC-2023-0071 \cite{rustsec2023_0071}, con plan de contingencia de regresar la firma a \texttt{ring} si el riesgo se volviera inaceptable (\rutarepo{docs/adr/0002-rsa-signing-crate-and-advisory.md}).
  \item \textbf{Autoridad de sellado demostrada.} El diseño delegó el sellado de tiempo en el sandbox de Cincel (\cref{sec:sellado-tiempo}); la vía demostrada de extremo a extremo es la TSA local con OpenSSL. El adaptador remoto está implementado y probado contra stub, pero la inestabilidad observada de la API y el costo del sandbox impidieron una campaña externa reproducible, como detalla \cref{sec:impl-spike-tsa}.
  \item \textbf{Dependencias del dominio.} El diseño declaró que el dominio dependería exclusivamente de \texttt{serde}, \texttt{thiserror}, \texttt{time} y \texttt{uuid} (\cref{sec:hexagonal-definicion}); la implementación añadió \texttt{zeroize}, aceptado porque es un crate utilitario sin E/S ni superficie de framework que codifica una propiedad de seguridad del dominio ---este valor se borra al liberarse--- y no un detalle de infraestructura (\rutarepo{docs/adr/0003-zeroize-secret-material-in-domain.md}).
  \item \textbf{Verificación de firmas X.509.} El diseño asumió el manejo de certificados con \texttt{x509-cert} y \texttt{der} (\cref{sec:inventario-crates}) sin anticipar que esas crates no verifican firmas; la verificación se implementó por re-codificación manual de la estructura to-be-signed (\rutarepo{docs/adr/0004-certificate-validation-implementation.md}), aceptando únicamente sha256WithRSAEncryption.
  \item \textbf{Escritura del paquete ZIP.} El diseño previó un paquete ZIP descargable sin fijar el mecanismo; se implementó un writer de entradas almacenadas que conserva bytes y evita un árbol adicional de dependencias. La actualización de MSRV elimina una restricción histórica, pero no exige sustituir el writer revisado (\rutarepo{docs/adr/0006-evidence-package-format.md}).
  \item \textbf{Superficie driving y persistencia.} El prototipo entrega CLI y API HTTP multiusuario; la interfaz Astro/Svelte permanece como andamiaje. PostgreSQL conserva usuarios, expedientes, pertenencias, documentos y auditoría con confirmación transaccional. Redis mantiene sesiones y controles efímeros. Los archivos locales se preservan para herramientas offline y migración; no son el almacén documental del servidor.
  \item \textbf{Sesiones opacas en lugar de JWT.} El diseño nombró JWT para la frontera autenticada. La implementación usa tokens aleatorios opacos porque la revocación inmediata ya exige consultar Redis y porque duplicar el rol dentro de un JWT permitiría que la autorización quedara obsoleta. Cada petición recarga estado activo y rol desde PostgreSQL; la sustitución permanece detrás del puerto \texttt{Session\allowbreak Store} (\rutarepo{docs/adr/0012-revocable-sessions-and-rbac.md}).
  \item \textbf{Clientes de datos síncronos.} El diseño eligió \texttt{sqlx} y \texttt{fred}; la implementación utiliza los clientes síncronos \texttt{postgres} y \texttt{redis}, con las versiones de seguridad compatibles con Rust 1.88. La frontera HTTP aísla los casos de uso en el pool bloqueante de Tokio. Antes de despliegue público deberán evaluarse pooling, TLS interno y clientes asíncronos.
  \item \textbf{Fitness functions.} De las cinco funciones de aptitud diseñadas (\cref{sec:fitness-functions}), la que exige un span de \texttt{tracing} por handler HTTP sigue diferida aunque los handlers ya están implementados; en su lugar se incorporó el umbral de cobertura del 90\,\% por crate criptográfico (FF-4, \cref{tab:impl-fitness}).
  \item \textbf{Alcance por recurso y funcionalidad pendiente.} La autorización combina rol, usuario activo, pertenencia vigente y asociación documental exacta. Cliente puede consultar metadatos asignados, pero su acceso documental sigue denegado por política. La persistencia de asignaciones no equivale a gestión completa de participantes procesales, audiencias, plazos o versiones documentales; esas funciones permanecen fuera del alcance implementado.
\end{itemize}

La separación por puertos permitió ampliar la persistencia y el acceso documental conservando la dirección de dependencias \cite{ford2023}.
